\title{Group Sequence Policy Optimization}

\begin{abstract}
We present Group Sequence Policy Optimization (GSPO), a robust and highly effective reinforcement learning method tailored for large language model training. Departing from conventional approaches that rely on token-wise importance ratios, GSPO operates by defining the importance ratio at the sequence level and applies sequence-wise clipping, reward assignment, and optimization procedures. Through empirical analysis, we show that GSPO not only outperforms the GRPO algorithm in training efficiency and overall performance but also enhances the stability of Mixture-of-Experts (MoE) reinforcement learning. Moreover, GSPO offers opportunities for streamlining RL infrastructure design. These strengths have driven significant advancements in the capabilities of recent Qwen3 models.
\end{abstract}

\section{Introduction}
\label{sec:intro}

Reinforcement learning (RL) has become a central methodology for enhancing the scalability of language models \citep{o1,dpsk-r1,qwq32b,qwen3}. Leveraging RL at scale enables language models to address challenging tasks—including advanced mathematics and programming—by facilitating more complex and extended reasoning.

For RL scaling to be effective as computational resources increase, it is crucial to ensure consistently stable and resilient training behavior. Nevertheless, prevailing RL techniques such as GRPO \citep{grpo} face significant stability challenges during the training of extremely large language models, frequently leading to drastic and unrecoverable failures in model performance \citep{qwen3, minimax-m1}. Such instability represents a significant obstacle to advancing the capabilities of language models via further reinforcement learning.

This work demonstrates that GRPO's instability originates from the flawed implementation and invalidation of importance sampling weights within its algorithm. Such deficiencies generate substantial training variance, which compounds as the generated response length increases; this effect is further intensified by the clipping procedure, leading eventually to the deterioration of the model.

To overcome these primary challenges, we introduce \textbf{Group Sequence Policy Optimization (GSPO)}, a novel reinforcement learning algorithm tailored for training large language models. GSPO's principal contribution is the rigorous formulation of the importance ratio using sequence likelihood \citep{click}, thereby conforming with the foundational concept of importance sampling. Moreover, GSPO calculates normalized rewards as the comparative advantages across multiple responses for each query, ensuring consistency between sequence-level reward attribution and the optimization objective.

Our experimental results reveal that GSPO markedly outperforms GRPO in terms of training consistency, resource utilization, and achieved outcomes. Notably, GSPO naturally addresses the stability issues encountered during RL training of extensive Mixture-of-Experts (MoE) models without requiring elaborate stabilization mechanisms, indicating its capacity to streamline RL system architecture. These advantages of GSPO were instrumental in achieving significant gains in the new Qwen3 models. We anticipate that GSPO will serve as a reliable and adaptable cornerstone for large-scale RL training processes involving language models, facilitating further progress in this area.

\section{Preliminaries}

\paragraph{Notation}
Throughout this work, we represent an autoregressive language model, which is parameterized by $\theta$, as a policy $\pi_\theta$. Queries are indicated by $x$, and the collection of all queries is represented by $\mathcal{D}$. For any response $y$ corresponding to a query $x$, the conditional probability under the policy $\pi_\theta$ is expressed as $\pi_\theta (y | x)=\prod_{t=1}^{|y|} \pi_\theta (y_t | x, y_{<t} )$, where $|y|$ refers to the total number of tokens in the response $y$. Each query-response tuple $(x, y)$ may be evaluated by a verifier $r$, producing a reward score $r(x, y) \in [ 0, 1 ]$.

\paragraph{Proximal Policy Optimization (PPO)} Using samples generated from the old policy $\pi_{\theta_\text{old}}$, PPO \citep{ppo} constrains the policy update within a proximal region of the old policy through the clipping mechanism. Specifically, PPO employs the following objective for policy optimization (we omit the KL regularization term hereinafter for brevity, as it is not the focus of this paper):

\begin{align}
\mathcal{J}_\text{PPO}(\theta) = \mathbb{E}_{ x \sim \mathcal{D},\, y \sim \pi_{\theta_\text{old}}( \cdot | x) }
\left[ \frac{1}{|y|} \sum_{t=1}^{|y|} 
\min \left( w_{t}(\theta) \widehat{A}_{t},  \, \mathrm{clip} \left( w_{t}(\theta), 1 - {\varepsilon}, 1 + {\varepsilon}\right) \widehat{A}_{t} \right)
\right],
\end{align}

Here, the token $y_t$ is assigned an importance ratio defined as $
w_{t}(\theta) = \frac{ \pi_{\theta} (y_{t} | x, y_{<t}) }{ \pi_{\theta_\text{old}} (y_{t} | x,y_{<t})} $. The advantage $\widehat{A}_{t}$ associated with $y_t$ is inferred by a separate value model, and the parameter $\varepsilon$ denotes the allowable clipping interval for importance ratios.

A primary obstacle to PPO's application stems from its substantial dependence on the value model. In typical implementations, the value model matches the policy model in scale, resulting in significant demands on both memory and computation. Additionally, PPO's algorithmic performance is fundamentally tied to the precision of the value estimates it generates. While the difficulty of constructing a trustworthy value model is well-known, expanding its reliability to accommodate longer outputs and more intricate tasks exacerbates the complexity of the issue.

\paragraph{Group Relative Policy Optimization (GRPO)} GRPO \citep{grpo} bypasses the need for the value model by computing the relative advantage of each response within a group of responses to the same query. Specifically, GRPO optimizes the following objective:

\begin{align}
\label{equ:grpo}
\mathcal{J}_\text{GRPO}(\theta) = \mathbb{E}_{ x \sim \mathcal{D},\, \{y_i\}_{i=1}^G \sim \pi_{\theta_\text{old}}( \cdot | x) }
\left[ \frac{1}{G} \sum_{i=1}^{G} \frac{1}{|y_i|} \sum_{t=1}^{|y_i|} 
\min \left( w_{i,t}(\theta) \widehat{A}_{i,t},  \, \mathrm{clip} \left( w_{i,t}(\theta), 1 - {\varepsilon}, 1 + {\varepsilon}\right) \widehat{A}_{i,t} \right)
\right],
\end{align}

where $G$ is the number of generated responses to each query $x$ (i.e., the group size), and the importance ratio $w_{i,t}(\theta)$ and advantage $\widehat{A}_{i,t}$ of token $y_{i,t}$ are:

\begin{align}
    w_{i,t}(\theta)=\frac{ \pi_{\theta} (y_{i,t} | x, y_{i,<t}) }{ \pi_{\theta_\text{old}} (y_{i,t} | x,y_{i,<t})},\quad\ 
    \widehat{A}_{i,t} = \widehat{A}_{i} = \frac{r(x, y_i) - \mathrm{mean} \left( \{ r(x, y_i) \}_{i=1}^G \right) }{ \mathrm{std} \left( \{ r(x, y_i) \}_{i=1}^G \right) },
\end{align}

respectively, where all the tokens in $y_i$ share the same advantage as $\widehat{A}_{i}$.

\section{Motivation}
\label{sec:motivation}

As models increase in scale, incorporate greater sparsity (such as in Mixture-of-Experts architectures), and generate longer responses, there is an increased requirement for large rollout batch sizes to fully leverage hardware capabilities during reinforcement learning. To enhance sample efficiency, a common approach involves subdividing the extensive rollout batch into several mini-batches, which are then used for gradient computation. This methodology inevitably results in an off-policy learning context: the responses $y$ originate from a previous policy $\pi_{\theta_\text{old}}$ instead of the present policy $\pi_\theta$ under optimization. As a consequence, mechanisms like clipping in PPO and GRPO become essential; they ensure that samples which deviate substantially from the current policy do not adversely affect the gradient estimation.

While mechanisms like clipping aim to manage this off-policy discrepancy, we identify a more fundamental issue in GRPO: \textit{its objective is ill-posed}. This problem becomes particularly acute when training large models on long-response tasks, leading to catastrophic model collapse. The ill-posed nature of the GRPO objective stems from a misapplication of importance sampling weights. The principle of importance sampling is to estimate the expectation of a function $f$ under a target distribution $\pi_\text{tar}$ by re-weighting samples drawn from a behavior distribution $\pi_\text{beh}$:

\begin{align}
\mathbb{E}_{z \sim \pi_\text{tar}} \left[ f(z) \right] = \mathbb{E}_{z \sim \pi_\text{beh}} \left[ \frac{ \pi_\text{tar} (z) }{ \pi_\text{beh} (z)} \, f(z) \right].
\end{align}

Importantly, the success of this method hinges on averaging across a large number of samples ($N \gg 1$) drawn from the behavior policy $\pi_\text{beh}$; only then does the importance weight $\frac{ \pi_\text{tar} (z) }{ \pi_\text{beh} (z)}$ adequately mitigate the distributional shift.

By contrast, GRPO leverages the importance weight $\frac{ \pi_{\theta} (y_{i,t} | x, y_{i,<t}) }{ \pi_{\theta_\text{old}} (y_{i,t} | x,y_{i,<t})}$ individually at each token position $t$. However, since each weight depends on only a single sampled token $y_{i,t}$ per token distribution $\pi_{\theta_\text{old}}( \cdot | x, y_{i, <t})$, it does not effectively correct for distribution mismatch. Rather, it introduces substantial variance into the gradient estimates during training. This effect accumulates over lengthy sequences, and is further intensified by the application of the clipping procedure. Our empirical findings indicate that this instability frequently results in catastrophic model collapse—a failure that, once triggered, is irrecoverable. Subsequent attempts to restore progress, whether by reverting to earlier model checkpoints, extensively adjusting hyperparameters like clipping boundaries, increasing generation length, or altering RL objectives, do not remedy the collapse.

These findings highlight an intrinsic flaw in the GRPO framework. Specifically, the ineffectiveness of the token-level importance weights underscores a fundamental principle: \textbf{the granularity of the optimization objective must correspond to that of the reward signal}. Given that the reward is assigned at the sequence level, imposing off-policy corrections at the level of individual tokens is misaligned. This observation leads us to abandon the token-level approach, instead focusing on computing importance weights and conducting optimization directly at the \textit{sequence level}.

\section{Algorithm}

\subsection{GSPO: Group Sequence Policy Optimization}

Although using the token-level importance weight $\frac{ \pi_{\theta} (y_{i,t} | x, y_{i,<t}) }{ \pi_{\theta_\text{old}} (y_{i,t} | x, y_{i,<t})}$ presents challenges within GRPO, we find that the \textit{sequence-level} importance weight $\frac{ \pi_{\theta} (y | x) }{ \pi_{\theta_\text{old}} (y | x)}$ carries a well-defined theoretical interpretation in the domain of language generation. Specifically, it quantifies the extent to which a response $y$ generated from $\pi_{\theta_\text{old}} (\cdot | x)$ differs from that under $\pi_{\theta} (\cdot | x)$. This measure aligns intuitively with the notion of sequence-level rewards, and offers a meaningful perspective on the behavior of the clipping mechanism.

Based on this straightforward observation, we propose the \textbf{Group Sequence Policy Optimization (GSPO)} algorithm. GSPO employs the following sequence-level optimization objective:

\begin{align}
\mathcal{J}_\text{GSPO} (\theta) =
\mathbb{E}_{ x \sim \mathcal{D},\, \{y_i\}_{i=1}^G \sim \pi_{\theta_\text{old}}( \cdot | x) }
\left[ 
\frac{1}{G} \sum_{i=1}^{G}
\min \left( s_{i}(\theta)  \widehat{A}_{i},  \, \mathrm{clip} \left( s_{i}(\theta), 1 - {\varepsilon}, 1 + {\varepsilon} \right) \widehat{A}_{i} \right) 
\right],
\label{equ:gspo}
\end{align}

where we adopt the group-based advantage estimation:

\begin{align}
\widehat{A}_{i} = \frac{r(x, y_i) - \mathrm{mean} \left( \{ r(x, y_i) \}_{i=1}^G \right) }{ \mathrm{std} \left( \{ r(x, y_i) \}_{i=1}^G \right) },
\end{align}

and define the importance ratio $s_{i}(\theta)$ based on sequence likelihood \citep{click}:

\begin{align}
s_{i}(\theta) = \left( \frac{ \pi_{\theta} (y_i | x) }{ \pi_{\theta_\text{old}} (y_i | x)} \right)^{\frac{1}{|y_i|}}
=
\exp \left( \frac{1}{|y_i|} \sum_{t=1}^{|y_i|} \log \frac{ \pi_{\theta} (y_{i,t} | x, y_{i,<t}) }{ \pi_{\theta_\text{old}} (y_{i,t} | x,y_{i,<t})} \right).
\end{align}

As a result, GSPO performs clipping at the response level, rather than at the token level, to eliminate highly ``off-policy'' samples from gradient calculations, which ensures alignment with both sequence-based reward allocation and optimization processes. Additionally, length normalization is incorporated in $s_{i}(\theta)$ to mitigate variance and maintain $s_{i}(\theta)$ within a consistent numerical range. Without this normalization, alterations in the likelihood of a small subset of tokens could cause substantial instability in the sequence-level importance ratio, and responses of different lengths would necessitate separate clipping ranges. Furthermore, it should be highlighted that the magnitude of clipping ranges utilized in GSPO tends to differ significantly from those in preceding methods (such as GRPO), attributable to the inherently distinct definitions of importance ratios in these frameworks.

\subsection{Gradient Analysis}
\label{subsec:gradient}

We can derive the gradient of the GSPO objective as follows (clipping is omitted for brevity):

\begin{align}
\nabla_{\theta} \mathcal{J}_\text{GSPO} (\theta)
=&\ 
\nabla_{\theta} \mathbb{E}_{ x \sim \mathcal{D},\, \{y_i\}_{i=1}^G \sim \pi_{\theta_\text{old}}( \cdot | x) }
\left[ 
\frac{1}{G} \sum_{i=1}^{G} s_{i}(\theta) \widehat{A}_{i}
\right] \\
= &\ 
\mathbb{E}_{ x \sim \mathcal{D},\, \{y_i\}_{i=1}^G \sim \pi_{\theta_\text{old}}( \cdot | x) }
\left[ 
\frac{1}{G} \sum_{i=1}^{G}
s_{i}(\theta) \widehat{A}_{i}
\cdot \nabla_{\theta} \log s_{i}(\theta)
\right] \\
=&\ 
\mathbb{E}_{ x \sim \mathcal{D},\, \{y_i\}_{i=1}^G \sim \pi_{\theta_\text{old}}( \cdot | x) }
\left[ 
\frac{1}{G} \sum_{i=1}^{G} 
\left( \frac{ \pi_{\theta} (y_i | x) }{ \pi_{\theta_\text{old}} (y_i | x)} \right)^{\frac{1}{|y_i|}} \widehat{A}_{i} 
\cdot \frac{1}{|y_i|} \sum_{t=1}^{|y_i|} \nabla_{\theta} \log \pi_{\theta} (y_{i,t} | x, y_{i,<t}) 
\right].
\label{equ:gspo_grad}
\end{align}

For comparison, the gradient of the GRPO objective is as follows (note that $\widehat{A}_{i,t} = \widehat{A}_{i}$):

\begin{align}
\nabla_{\theta} \mathcal{J}_\text{GRPO} (\theta) 
=&\ 
\nabla_{\theta} \mathbb{E}_{ x \sim \mathcal{D},\, \{y_i\}_{i=1}^G \sim \pi_{\theta_\text{old}}( \cdot | x) }
\left[ \frac{1}{G} \sum_{i=1}^{G} \frac{1}{|y_i|} \sum_{t=1}^{|y_i|} 
w_{i,t}(\theta) \widehat{A}_{i,t}
\right] \\
=&\
\mathbb{E}_{ x \sim \mathcal{D},\, \{y_i\}_{i=1}^G \sim \pi_{\theta_\text{old}}( \cdot | x) }
\left[ \frac{1}{G} \sum_{i=1}^{G} \widehat{A}_{i} 
\cdot \frac{1}{|y_i|} \sum_{t=1}^{|y_i|} 
\frac{ \pi_{\theta} (y_{i,t} | x, y_{i,<t}) }{ \pi_{\theta_\text{old}} (y_{i,t} | x,y_{i,<t})} 
\nabla_{\theta} \log \pi_{\theta} (y_{i,t} | x, y_{i,<t})  
\right].
\end{align}

The primary difference between GSPO and GRPO centers on \textit{the method by which they assign weights to the gradients of token log likelihoods}. GRPO determines these weights using the so-called "importance weight" $\frac{ \pi_{\theta} (y_{i,t} | x, y_{i,<t}) }{ \pi_{\theta_\text{old}} (y_{i,t} | x,y_{i,<t})}$ for each token. These varying weights—falling within $(0, 1+\varepsilon]$ when $\widehat{A}_{i} >0$ or $[1-\varepsilon, +\infty)$ when $\widehat{A}_{i} < 0$—are substantial and cannot be ignored, as their cumulative effects may introduce instability and unpredictable behavior in the training process. GSPO, by contrast, applies uniform weighting to every token in a sequence, thus sidestepping the instability issues inherent in GRPO.

\subsection{GSPO-token: A Token-level Objective Variant}

In scenarios like multi-turn RL, we may desire a finer-grained advantage adjustment than the sequence level. To this end, we introduce a token-level objective variant of GSPO, namely \textbf{GSPO-token}, to allow token-wise advantage customization:

\begin{align}
\mathcal{J}_\text{GSPO-token}(\theta)
=
\mathbb{E}_{ x \sim \mathcal{D},\, \{y_i\}_{i=1}^G \sim \pi_{\theta_\text{old}}( \cdot | x) }
\left[ 
\frac{1}{G} \sum_{i=1}^{G} \frac{1}{|y_i|} \sum_{t=1}^{|y_i|} 
\min \left( s_{i,t}(\theta) \widehat{A}_{i,t},  \, \mathrm{clip} \left( s_{i,t}(\theta), 1 - {\varepsilon}, 1 + {\varepsilon} \right) \widehat{A}_{i,t} \right) 
\right],
\label{equ:gspo-token}
\end{align}

where

\begin{align}
s_{i,t}(\theta) 
= \mathrm{sg} \left[ s_{i}(\theta) \right]  \cdot \frac{ \pi_{\theta} (y_{i,t} | x, y_{i,<t}) }{ \mathrm{sg} \left[ \pi_{\theta} (y_{i,t} | x, y_{i,<t}) \right] },
\end{align}

and $\mathrm{sg}[\cdot]$ denotes only taking the numerical value but stopping the gradient, corresponding to the \texttt{detach} operation in PyTorch. The gradient of GSPO-token can be derived as:

\begin{align}
\nabla_{\theta} \mathcal{J}_\text{GSPO-token}(\theta)
=&\ 
\nabla_{\theta} \mathbb{E}_{ x \sim \mathcal{D},\, \{y_i\}_{i=1}^G \sim \pi_{\theta_\text{old}}( \cdot | x) }
\left[ 
\frac{1}{G} \sum_{i=1}^{G}
\frac{1}{|y_i|} \sum_{t=1}^{|y_i|} 
s_{i,t}(\theta) \widehat{A}_{i,t}
\right] \\
=&\ 
\mathbb{E}_{ x \sim \mathcal{D},\, \{y_i\}_{i=1}^G \sim \pi_{\theta_\text{old}}( \cdot | x) }
\left[ 
\frac{1}{G} \sum_{i=1}^{G} s_i (\theta) \cdot \frac{1}{|y_i|} \sum_{t=1}^{|y_i|} 
\widehat{A}_{i,t} \frac{ \nabla_{\theta} \pi_{\theta} (y_{i,t} | x, y_{i,<t}) }{ \pi_{\theta} (y_{i,t} | x, y_{i,<t}) }
\right] \\
=&\ 
\mathbb{E}_{ x \sim \mathcal{D},\, \{y_i\}_{i=1}^G \sim \pi_{\theta_\text{old}}( \cdot | x) }
\left[ 
\frac{1}{G} \sum_{i=1}^{G}  \left( \frac{ \pi_{\theta} (y_i | x) }{ \pi_{\theta_\text{old}} (y_i | x)} \right)^{\frac{1}{|y_i|}}
\cdot \frac{1}{|y_i|} \sum_{t=1}^{|y_i|}
\widehat{A}_{i,t} \nabla_{\theta} \log \pi_{\theta} (y_{i,t} | x, y_{i,<t})  
\right].
\label{equ:gspo-token_grad}
\end{align}

Observe that $\frac{ \pi_{\theta} (y_{i,t} | x, y_{i,<t}) }{ \mathrm{sg} \left[ \pi_{\theta} (y_{i,t} | x, y_{i,<t}) \right] }$ evaluates to 1, thereby making $s_{i,t}(\theta)$ equivalent in numerical value to $s_{i}(\theta)$. By examining Equation~\eqref{equ:gspo} alongside~\eqref{equ:gspo-token}, as well as Equation~\eqref{equ:gspo_grad} with~\eqref{equ:gspo-token_grad}, it is evident that GSPO-token and GSPO yield identical results in terms of their optimization objectives, clipping constraints, and theoretical gradients whenever all token advantages within the response $y_i$ are set uniformly (i.e., $\widehat{A}_{i,t} = \widehat{A}_{i}$). Nevertheless, GSPO-token provides greater flexibility, permitting the adjustment of token advantages individually.

\section{Experiments and Discussion}

\subsection{Empirical Results}

We conduct experiments utilizing a cold-start model that has been fine-tuned from Qwen3-30B-A3B-Base, and present both the training reward trajectories and evaluation metrics on several benchmarks: AIME'24 (reporting average Pass@1 across 32 samples), LiveCodeBench (spanning 202410 to 202502, reporting average Pass@1 over 8 samples), and CodeForces (measured by Elo Rating). Throughout the RL training phase, each batch of rollout data is subdivided into four mini-batches for applying gradient updates. For GSPO, we specify the lower and upper clipping thresholds in Equation~\eqref{equ:gspo} as 3e-4 and 4e-4, respectively. In our comparative analysis with the GRPO baseline, the lower and upper clipping parameters in Equation~\eqref{equ:grpo} are adjusted to 0.2 and 0.27, with these values carefully selected to maintain fairness in our experiments. It is important to highlight that GRPO requires the Routing Replay training scheme for stable MoE RL convergence, a detail we elaborate on in \S~\ref{subsec:routing_replay}. By contrast, \textit{GSPO eliminates the necessity for this procedure}.

As illustrated in Figure~\ref{fig:results}, GSPO-based training maintains consistent stability across its progression. Our analysis indicates that \textbf{GSPO consistently enhances performance by scaling training compute, frequently refreshing the query set, and lengthening generation output}. Additionally, GSPO exhibits greater training efficiency relative to GRPO, attaining superior accuracy and benchmark scores with equivalent training resources and query usage. Lastly, GSPO has been effectively employed in the RL training of the most recent Qwen3 models, which provides compelling evidence for its capability in harnessing RL scaling to advance large language models.

\subsection{Curious Observation on Clipping Fractions}

GSPO distinguishes itself from GRPO primarily through its approach to clipping: rather than applying clipping at the token level, GSPO clips entire response sequences. Figure~\ref{fig:clipping} demonstrates that GSPO involves clipping a markedly greater proportion of tokens—by nearly two orders of magnitude—relative to GRPO, irrespective of adjustments made to the clipping intervals. Notably, even with the substantial increase in clipped tokens and the resultant reduction in tokens available for model training or gradient computation, GSPO demonstrates superior training efficiency compared to GRPO. This seemingly paradoxical outcome—that training efficiency improves even as a larger share of tokens are clipped—suggests that GRPO's gradient estimations at the token level are subject to considerable noise and lack efficacy for extracting value from samples. On the other hand, GSPO's methodology of generating gradients at the sequence level yields a learning signal that is both more robust and more conducive to efficient training.

\subsection{Benefit of GSPO for MoE Training}
\label{subsec:routing_replay}

\paragraph{Background}
In contrast to the reinforcement learning (RL) optimization of dense neural networks, Mixture of Experts (MoE) models present distinct stability issues due to their inherently sparse pattern of activation. Notably, our observations with the GRPO algorithm reveal that MoE models suffer from substantial \textit{expert-activation volatility}, which can impede RL training from achieving stable convergence. Specifically, following one or more steps of gradient updates, the subset of experts activated for generating an identical response may shift appreciably. Taking the 48-layer Qwen3-30B-A3B-Base model as an example, after each RL gradient step, analysis of the same sampled rollout shows that approximately 10\% of the active experts change when comparing activations under the newly updated policy $\pi_{\theta}$ to those under the preceding policy $\pi_{\theta_\text{old}}$.

This variability, which is especially acute in deeper MoE architectures, causes the token-level importance weights $w_{i,t}(\theta) = \frac{ \pi_{\theta} (y_{i,t} | x, y_{i,<t}) }{ \pi_{\theta_\text{old}} (y_{i,t} | x,y_{i,<t})}$ to become highly unstable, thereby undermining their reliability as noted in \S~\ref{sec:motivation} and~\ref{subsec:gradient}. As a result, this instability obstructs the steady progress of RL training.

\paragraph{Our Previous Approach} In our earlier work, we introduced the \textbf{Routing Replay} training method to address this issue. Our solution involves recording which experts were activated under $\pi_{\theta_\text{old}}$ and then reproducing these same routing decisions for $\pi_{\theta}$ during the calculation of the importance ratios $w_{i,t}(\theta) = \frac{ \pi_{\theta} (y_{i,t} | x, y_{i,<t}) }{ \pi_{\theta_\text{old}} (y_{i,t} | x,y_{i,<t})}$. This procedure guarantees that, for each token $y_{i,t}$, both $\pi_{\theta} (y_{i,t} | x, y_{i,<t})$ and $\pi_{\theta_\text{old}} (y_{i,t} | x,y_{i,<t})$ utilize an identical set of activated experts, thus maintaining consistency in the token-level importance ratios and supporting updates to a stable network architecture during optimization steps. As illustrated in Figure~\ref{fig:routing_replay}, Routing Replay plays a pivotal role in ensuring the GRPO training of MoE models converges as expected.

\paragraph{Benefit of GSPO} While Routing Replay allows GRPO training for MoE models to converge effectively, it introduces extra memory and communication costs and restricts the model’s usable capacity by necessitating repeated routing modes. In comparison, GSPO—demonstrated in Figure~\ref{fig:results}—does away with the reliance on Routing Replay. It permits direct computation of importance ratios $s_i(\theta)$, ensuring proper convergence and stable optimization. The central observation is that GSPO operates solely based on the sequence likelihood ($\pi_{\theta} (y_i | x)$), without dependence on the likelihood of individual tokens ($\pi_{\theta} (y_{i,t} | x, y_{i,<t})$). Because the MoE architecture consistently preserves robust language modeling, the sequence likelihood remains steady. Thus, GSPO effectively addresses and eliminates the expert-activation instability inherent in MoE models, removing the need for elaborate fixes such as Routing Replay. This results in a more streamlined and reliable training process and permits unrestricted utilization of the model’s full expert capacity.

\subsection{Benefit of GSPO for RL Infrastructure}

Due to differences in numerical precision between training frameworks such as Megatron and inference systems like SGLang and vLLM, the standard approach has been to recalculate sampled response likelihoods under the previous policy $\pi_{\theta_\text{old}}$ using the training engine. Nonetheless, since GSPO bases its optimization on sequence-level likelihoods—as opposed to the more sensitive token-level metrics—it exhibits a greater robustness to such precision variations. As a result, GSPO enables direct utilization of likelihood outputs from the inference engine for optimization purposes, eliminating the necessity for additional recomputation with the training system. This property is particularly advantageous in contexts involving partial rollouts, multi-turn reinforcement learning, and frameworks that decouple training from inference.

\section{Conclusion}

We introduce Group Sequence Policy Optimization (GSPO), a novel algorithm for reinforcement learning aimed at training large language models. GSPO employs sequence-level importance ratios derived from sequence probabilities, utilizing importance sampling as its guiding principle. It applies clipping, rewarding, and optimization procedures at the sequence level. Empirically, GSPO outperforms GRPO in terms of training stability, efficiency, and overall performance, with marked advantages for large-scale RL in MoE model training. This contributes significantly to the substantial advancements observed in recent Qwen3 models. As a robust and scalable algorithmic platform, GSPO enables ongoing expansion of RL and heralds foundational progress in the development of intelligent systems.

\end{document}